\documentclass[11pt]{article}
\usepackage[margin=1in]{geometry}
\usepackage{amsmath,amssymb,booktabs,hyperref}
\title{$1\times1 = 2$ and $0/0 = 1$:\\[4pt]\large A Counting System in Which No Unit Is Lost, with a Fold Notation for Numbers}
\author{Turaq Rayn-McLean\\Independent researcher}
\date{November 27, 2026 (version 5.8)}
\emergencystretch=3em
\begin{document}
\maketitle
\begin{center}\emph{For Terence McKenna.}\end{center}

\begin{abstract}
We define a counting system in which a unit is never lost in multiplication. Ordinary products hold wherever a factor already contains the unit; where bare units meet bare units, every unit present is counted, so $1\times1 = 2$ and the unit multiplied by itself $n$ times is $n$. The system separates two zeros (nothingness and absence) and two ones (the unit and the literal one). With absence as the zero, $0/0 = 1$: an empty slot measured by an empty slot is one unit of room. Division by a bare unit may hold two values at once, and the distributive law is not part of the system. Every number has a count root and a measure root; the order-$n$ count root of $n$ is $1$ for every $n$, while the ordinary $n^{1/n}$ reaches $1$ only at infinity. We also give a fold notation that writes $0$, $1$, $2$, $3$ as $000$, $010$, $101$, $111$ (wing, fold, wing), builds every number by halving into a mirror-symmetric string, and expands it in orders by one rewrite ($0\to000$, $1\to010$) that keeps every unit and its place while the zeros carry the scale. A second, divisor build writes $n = d\times m$ as $d$ copies of $m$ with a zero between copies that grows one order per level; it keeps exactly $n$ units, gives a prime a single face and a divisor-rich number many. We prove (in Lean 4 with Mathlib) that three is the only word length above one in which every unit count has exactly one mirror word, which is why the notation is triaxial; a ternary mirror build then reproduces the four words and the orders by itself, keeps exactly $n$ units, and has no drift, while the binary halving build drifts by exactly $5/4$ per octave. We also prove that a number and its digit mirror share their remainder by $9$ and, up to a sign set by the digit count, by $11$, and that $2$ and $5$ act alike through $10$: when the decimal period of $1/p$ is odd their orders mod $p$ carry the same power of $2$, and when it is even they are never both of odd order. The system agrees with ordinary arithmetic on every expression that contains a number holding the unit, and departs from it only where bare units meet. Version 4 adds facing marks, an equal-width zoom in which the comma count is the scale, and a ladder shorthand that writes any number by one pair of neighbours per level. For the number of units $u(n)$ in the halving build we prove in Lean that $u(n) = n$ exactly at $2^a$ and $3\cdot2^a$, that $u(n)$ is never $2n$, and that the odd numbers below two units per number are exactly the two strands $2^a+1$ and $3\cdot2^a+1$ and the three points $1$, $11$, $15$. Version 5 adds one code with three faces, the reading of halving as a number's neighbours and of turns as the number itself, Lean proofs for every set of turns (no number holds twice itself in units, and every exception sits beside a self-counting number), the relationship by which each split between consecutive primes hands its small mirrors to its center and half gap, the hop of walls and seat steps, the hierarchy of seat rules, even primes, offsets of eleven, nine and cubes, a construction that names every prime in turn up to $10^7$, two as the zero of the primes, and the wall law: a prime block $r = 6k_r+f_r$ closes the wing $f$ of wall $k$ exactly when $k\equiv f f_r k_r \pmod r$ (proven in Lean), the two-point law $u(c2^a+j) = u(c)v(2^a+j) - d(c)v(j)$, and an audit placing every claim as proven ($77$ of $141$, in one Lean library of $330$ theorems), checked, measured, definition or reading. Version 5.3 runs the wall law from $2$ and $3$ alone through every seat to $4.295\times10^9$ and to the last seats below $2^{64}$, exact throughout. Version 5.4 proves the wall walk exact at every size and raises the proven claims to $91$ of $141$, in one Lean library of $398$ theorems, with laws found by PySR and then proven. Version 5.5 adds one formula for the primes in wall coordinates, proven, and says plainly that it is a computation ($92$ of $142$ claims proven, $404$ theorems). Version 5.6 states the work as a solve of the prime pattern as a structure, proves the wall law at every width and the walk on walls of $30$, proves the wall count equal to Mathlib's prime count, proves one hand-off for every prime block, and walks to $10^{13}$ in $1051$ seconds ($93$ of $142$ claims proven, $422$ theorems). Version 5.7 ties the wall count to the Riemann hypothesis in Lean: from the proved hypothesis, the wall count is within $C N^{1-\delta}$ of $\operatorname{Li}(N)$. Version 5.8 sharpens this to $C\sqrt{N}(\log N)^2$, and to $C N^{1/2+\delta}$ for every $\delta>0$. Reference code and the Lean proofs are attached.
\end{abstract}

\section*{What this is, and what it is not}
This paper gives a counting system, a notation for numbers built from mirrors, folds and turns, and what that notation shows about the primes. Its firm results are theorems proved in Lean~4 with Mathlib (listed with their projects in the code bundle) about a new sequence, the unit count $u$, and about the structure of every set of turns. Its readings of the primes are relationships and families: how a split between consecutive primes hands its small mirrors to its center and half gap, how seat rules stack, how the range for the next prime closes as maps are added, and how a construction from passed numbers names every prime in turn to $10^7$. The author presents it as a solve of the prime pattern as a structure: where every prime sits, the laws between consecutive primes, and a walk that names each one exactly, proven in Lean for every size. It is not a formula that skips the computation, and the qQq remains. It is a common relationship with its families, by a novel method, proven where Lean could reach and measured where it could not. It does not claim one formula for every prime; the author's view is that there is none, that the primes are a relationship with families, and that a residue (the qQq) remains at every scale. Where a fact matches earlier work it is noted, with thanks, as the same fact reached by a different route; the method and its application are the author's.

\section*{Terms and symbols}
The terms and symbols below are the author's and mean what they mean in this system. The system is the author's own version of mathematics; it is novel, and it is not meant to be read as a translation into, or out of, other systems.
\begin{description}\setlength{\itemsep}{1pt}
\item[$0_n$, $0_a$] Nothingness, and absence (an empty seat a unit could fill).
\item[$1_u$, $1_U$] The unit, and the unit that holds anything in the package, including infinity.
\item[$000, 010, 101, 111$] The words for $0, 1, 2, 3$: wing, centre, wing.
\item[wing, centre (center), hinge] The two outer blocks of a picture, its middle block, and a zero centre.
\item[$w,c,w$] One code: $w$ is half of $n$ (rounded down), $c$ is $0$ (hinge) or $1$ (the unit).
\item[comma] One level down into the notation; the number of commas is the number of levels.
\item[$|$] The mirror seat: fold back on the digit it touches; between commas, a fold at that level.
\item[$a,|b|,c$] A product face: $a\times b\times c$ (with bare units, $1,|1|,1 = 3$).
\item[$<\,>$] Facing: an arrow pointing in marks an operator acting on the centre; pointing out, an observer.
\item[$\text{\textasciicircum}$, v] One scale step up (out, the telescope), one step down (in, the microscope).
\item[$|\text{\textasciicircum}x$] $x$ mirrored across one scale step, both kept: $x$ times (base $+1$); $|\text{\textasciicircum}1$ is $3$ in base $2$ and $11$ in base $10$.
\item[$a,b,a\!:\!s$] A picture read at scale $s$: centre plus $(s-1)$ wings ($s=1$ the centre alone, $s=2$ halving, $s=3$ the turn).
\item[reflection, turn] The halving mirror (two-based) and the thirding mirror (three equal copies).
\item[faces] The forms of one number in one code: add, copies, product, reflection.
\item[s, m, lm] The standard line, the mirror line (offset $11$: $n_m = 11 - n_s$; and $22$), and the lake mirror (upside down, going down into scale).
\item[strand] A family of numbers laid by one rule with a period or a law, running as far as it runs.
\item[seat, open seat, missing wing] A place a number can sit; a seat no strand reaches; a wing whose seat is closed.
\item[left-out] An even number that folds into equal copies of an even block other than $2$.
\item[even prime] A prime place on the even line (place $k$ at $2k$); also the evens not built from two evens.
\item[hop, wall] The step from one prime to the next; a multiple of $6$ crossed on the way.
\item[qQq] What is not in the finite map: whatever lies outside the map at the time, unknown as it is; the mirror boundary that sends the image back. For the primes, the residue that remains outside every map laid so far. It is never expected to vanish, and each added map moves it.
\item[frayed rope] Strands that run far and strands that cut short; a strand ending is information, not failure.
\end{description}

\section{Two zeros and two ones}
Ordinary arithmetic has one zero and one one. This system has two of each.
\begin{center}\begin{tabular}{lll}\toprule
Symbol & Name & Meaning\\\midrule
$0_n$ & nothingness & nothingness itself; the infinite absence of absence\\
$0_a$ & absence & the place for a unit that is not present but could be\\
$1_u$ & the unit & any unit: one apple, one dozen, one carton\\
$1_c$ & literal one & the count one\\\bottomrule
\end{tabular}\end{center}
A carton of twelve eggs is $1_u$ of eggs. ``I am out of apples; a box comes tomorrow'' is $0_a$: an absence shaped by the unit it is missing. Nothingness has no unit to fill it.

\section{The counting rule}
\paragraph{Containment.} A number $n>1$ is $1,2,\dots,n$, so it already holds the unit. Hence $1\times n = n$; the unit does not carry.

\paragraph{Bare units.} When every factor is a bare unit, none holds another, and every unit present is counted:
\[ 1\times1 = 2,\qquad 1\times1\times1 = 3,\qquad \underbrace{1\times\cdots\times1}_{n} = n. \]
This holds for $1_u$ and $1_c$ alike: one apple times one more apple is two apples. With $u$ one dozen, one dozen times one dozen is two dozen, the way cartons are counted in practice.

\paragraph{Whole expression.} The rule reads the whole expression at once. If every factor is a bare one, the result is the number of ones. If any factor holds the unit, the ones drop out and the remaining factors multiply as usual:
\[ 1\times1\times3 = 3,\qquad 1\times2 = 2,\qquad 2\times3 = 6. \]
The number $2$ is a count $(1,2)$, not the product $1\times1$ handed forward, so grouping does not change a result. The expression $1\times(1+1)$ reads ``one, then the second one'', that is $1,2$, and gives $1,2,3 = 3$, the same as $1\times(1\times1)$. Among bare units the product counts the way a sum does; elsewhere it is ordinary multiplication. In one sentence: multiplying accounts for every unit present.

\paragraph{No distribution.} Multiplication and addition are separate acts, and the distributive law is not part of the system. Leaving out a law of ordinary algebra is a known move: near-rings keep only one side of distribution, and the octonions are not associative.

\paragraph{Agreement.} Wherever a number holding the unit is present, the system agrees with ordinary arithmetic. It departs only where bare units meet, which is exactly where ordinary arithmetic sets $1\times1 = 1$ and a unit is lost.

\section{Division, and $0/0 = 1$}
Dividing by the unit leaves what you had, as a unit:
\[ 1/1 = 1,\qquad \infty/\infty = \infty,\qquad 0_n/1_u = \text{a nothingness unit},\qquad 0_a/1_u = \text{a unit of absence}. \]
Division by a bare unit may hold two values at once: $2/1$ asks what times $1$ gives $2$, and both $1$ (since $1\times1 = 2$) and $2$ (since $1\times2 = 2$) answer. The two zeros give $0/0$ two readings:
\[ 0_a/0_a = 1,\qquad 0_n/0_n = 0_n. \]
Ordinary arithmetic leaves $0/0$ undefined because it has one zero and cannot say which is meant. With absence, an empty slot measured by an empty slot is one unit of room. With nothingness, nothing is there to measure.

Powers follow the same count: $1^n = n$, so $1^0$ is zero bare units, which is $0_n$; and the unit multiplied by itself without end is without end.

\section{Count and measure}
Every number has two roots. Since the unit multiplied by itself $n$ times is $n$, the order-$n$ count root of $n$ is $1$ for every $n$: $\sqrt2 = 1$, $\sqrt[3]3 = 1$, $\sqrt[4]4 = 1$. The measure root is the ordinary $n^{1/n}$:
\begin{center}\begin{tabular}{rl@{\qquad}rl}\toprule
$n$ & $n^{1/n}$ & $n$ & $n^{1/n}$\\\midrule
2 & 1.4142135623730951 & 6 & 1.3480061545972777\\
3 & 1.4422495703074083 & 7 & 1.3204692477561237\\
4 & 1.4142135623730951 & 8 & 1.2968395546510096\\
5 & 1.379729661461215 & 10 & 1.2589254117941673\\\bottomrule
\end{tabular}\end{center}
The measure root reaches $1$ only as $n\to\infty$; the count root is $1$ at every step. The two disagree at every finite $n$ and meet only at infinity. They are furthest apart at $3$, the whole number nearest the peak at $e$ ($e^{1/e} = 1.444667861009766$). The count side is rational throughout; the measure side keeps its irrationals. The system does not remove $\sqrt2 = 1.4142135623730951$; it sets a count reading beside it.

\paragraph{Two yardsticks.} If a $1\times1$ square counts as $2$, a square of side $s$ holds $2s^2$, and the side of a square holding $n$ is $\sqrt{n/2}$. So $\sqrt3$ reads both $\sqrt3/1 = 1.7320508075688772$ and $\sqrt3/\sqrt2 = 1.224744871391589$. On the unit cube these are the body diagonal measured in edges and in face diagonals; the reciprocal $\sqrt{2/3} = 0.816496580927726$ is the body-to-face projection, and the square-yardstick root of $1$ is $1/\sqrt2 = 0.7071067811865475$. ``Three is $1,1,1$'' is the far corner of the unit cube, one step on each axis, and its cube root is $1$.

\section{Fold notation}
Each number is written as wing, fold, wing. The four words are
\[ 0 = 000,\qquad 1 = (0,1,0) = 010,\qquad 2 = (1,0,1) = 101,\qquad 3 = (1,1,1) = 111. \]
Zero is empty throughout. One is whole and has no wings. Two has an empty fold and one unit on each wing: $1\times1 = 2$ is the unit meeting itself and moving out to both wings. Three is the full line, both wings and the fold. These codes are written forms, not sums. Every number is built by halving:
\[ B(2k)=B(k),\,0,\,B(k),\qquad B(2k{+}1)=B(k),\,B(k{+}1),\,B(k), \]
with $B(1), B(2), B(3)$ the words above and $0$ the zero word. For example $4 = 101,000,101$; $5 = 2,3,2 = 101,111,101$; $6 = 111,000,111$; and
\[ \scriptstyle 17 = 8,9,8 = 101,000,101,000,101,000,101,\;101,000,101,101,111,101,101,000,101,\;101,000,101,000,101,000,101. \]
Every build reads the same from either wing and has a single centre triple, the hub. An even number's hub is the zero between its two halves; an odd number inherits the hub of its middle block. The number of triples obeys $T(1)=T(2)=T(3)=1$, $T(2k)=2T(k)+1$, $T(2k+1)=2T(k)+T(k+1)$.

\section{Orders}
The notation expands without end by one rewrite, applied to every digit at once:
\[ 0 \to 000,\qquad 1 \to 010. \]
Applying it once gives the words of order 1:
\[ 0 = 000,000,000,\quad 1 = 000,010,000,\quad 2 = 010,000,010,\quad 3 = 010,010,010, \]
and again gives order 2, and so on. The rewrite can be applied to the words before building or to the finished build; the result is the same string (checked for every $n<300$ at orders $0$, $1$, $2$), so an order is one step taken by the whole number at once. Mixing orders inside one number (writing some symbols at one order and others at another) changes the counts and is not allowed.

Across orders the units stay fixed: every number keeps the same number of ones, a one at digit $i$ moves to digit $3i+1$, and the centre stays the centre (proven, \texttt{one\_moves}, \texttt{centre\_stays}). What grows is the length, three times per order, and the zeros: the left wing of $11$ holds $10$ ones at every order and $3$, $30$, $111$ zeros at orders $0$, $1$, $2$. Each order is a step down into scale in which the units are held and the emptiness grows. Statements about the units, their places and the hub are true at every order; zero counts and distances are readings at a stated order. In the limit one becomes a single point at the centre of its cell, whole and without wings. (The word $101$ has the shape of Cantor's middle-third rule, which keeps the two outer thirds; $010$ is its mirror, which keeps only the centre.)

\section{Why three: the triaxial theorem}
A mirror word reads the same from both ends. For a word length $L$, count the mirror words with $w$ units. For $L=2$ the counts by $w = 0,1,2$ are $1,0,1$: a lone unit cannot sit in a two-slot mirror. For $L=3$ they are $1,1,1,1$; for $L=4$, $1,0,2,0,1$; for $L=5$, $1,1,2,2,1,1$.

\textbf{Theorem (proved in Lean 4 with Mathlib).} Every unit count $0,\dots,L$ has exactly one mirror word of length $L$ if and only if $L \in \{0, 1, 3\}$. For every $L\ge4$ the two words $1\,0\cdots0\,1$ and $0\,1\,0\cdots0\,1\,0$ both carry two units.

So three is the only width at which a mirror code is clean, and the four words $000, 010, 101, 111$ are exactly its mirror words: the notation is binary on three axes.

\section{The ternary build}
Split by three instead of two: write $n = 3k + r$ as three copies,
\[ r=0:\ k,k,k \qquad r=1:\ k,\,k{+}1,\,k \qquad r=2:\ k{+}1,\,k,\,k{+}1. \]
At level one this gives $000, 010, 101, 111$; at level two it gives exactly the order-1 words $000,000,000$; $000,010,000$; $010,000,010$; $010,010,010$. Every string is a mirror, each number has its own string, and the number of units equals $n$ exactly, with no excess (checked for $n\le 729$). Building one level deeper is the same as one order of the rewrite $0\to000$, $1\to010$. A centre unit appears exactly when $n$ is odd.

\textbf{The binary drift.} The halving build writes an odd number $2k+1$ as $k,(k{+}1),k$, both wings in full, which is a ternary split of $3k+1$ units where only $2k+1$ exist; it is the only source of excess. Over the octave $[2^k, 2^{k+1})$ the binary build holds $(17\cdot5^{k-1} - 2^k)/3$ units, growing $\times5$ against $\times4$ for the numbers: a drift of exactly $5/4$ per octave. The ternary build has none. The numbers with no excess in the binary build are exactly $2^k$ and $3\cdot2^k$, and the ternary build has no excess at any $n$ (both proven in Lean; an earlier version listed $3^k$ and $2\cdot3^k$ for the ternary build in error).

\section{Mirrors and remainders}
\textbf{Theorem (proved in Lean).} For any digit list, a number and its digit mirror leave the same remainder by $9$; by $11$ they leave the same remainder when the digit count is odd and opposite remainders when it is even. Hence if one side is a prime of three or more digits, the other side is never divisible by $3$ or $11$.

\textbf{Theorem (proved in Lean): $2$ and $5$ through $10$.} Let $p$ be a prime not dividing $10$. If the orders of $2$ and $5$ mod $p$ are both odd, so is the order of $10$; so when the decimal period of $1/p$ is even, $2$ and $5$ are never both of odd order. If the order of $10$ is odd, the orders of $2$ and $5$ carry the same power of $2$. The shapes $2$ and $5$ also swap under both the top-bottom and the left-right mirror of the seven-segment digits.

\textbf{Theorem (proved in Lean): the carry level.} For integers $x<y$, among $x{+}1,\dots,y$ exactly one integer has the largest power of $2$, and that power is $2^{\lfloor\log_2(x \oplus y)\rfloor}$. This is the level at which the halving build of $y$ first shares a block with that of $x$.

\section{The divisor build}
The halving build writes every number the same way. A second build writes a number through its factors. For $n = d\times m$, write $m$ a total of $d$ times, with a zero between neighbouring copies. The leaves are the words of $1$, $2$, $3$; a prime $p>3$ cannot be split into equal copies and is written as $p$ units, $010$ joined by $000$. The gaps between copies sit above that: the innermost gaps are $000,000,000$ (one order above the zero inside a prime), the next level out $27$ zeros, the next $81$, and so on, one order per level. For example
\[ 12 = 2\times2\times3 = 111,\,000,000,000,\,111,\;(27\text{ zeros}),\;111,\,000,000,000,\,111, \]
which is $6,0,6$ with the higher-order zero in the middle. Every arrangement of the prime factors is kept as one face of the number: $12$ also has the faces for the factor orders $2,3,2$ and $3,2,2$. In every face the number of ones is exactly $n$, every face reads the same from either end, every arrangement gives its own string, and no composite face reads like a prime (all checked for $n<2000$). A prime has one face; a number rich in divisors has many. The halving build treats every number alike; the divisor build is where primes and divisor-rich numbers part.

If the zero between copies were the same at every level, the levels would flatten (24 would read only as eight copies of $3$); growing the zero by one order per level is what keeps each level visible. The first gap must also sit one order above the zero inside a prime: otherwise a number such as $35 = 5\times7$, whose only gap is the first, reads exactly as $35$ written as a prime (the error of version 2.0).

\section{Facing marks and the zoom}
The halving build is kept, with two changes that make it one rule for every number. First, all blocks at one level are written at the same width: a narrower block is zoomed in by the rewrite $0\to000$, $1\to010$ until it matches, so every width is a power of $3$; the width is $3^{k+2}$ exactly for $3\cdot2^k < n \le 3\cdot2^{k+1}$ (proven, \texttt{width\_closed}). In $14 = 7,0,7$ with $7 = 3,4,3$, the $3$ is written $010,010,010$ beside the nine-digit $4$. Second, a mark $|$ shows where the mirror sits. It is a facing, not a sign: it says ``read back from here''. A mark beside a digit folds back on that digit, which is shared and written once, so $10| = 101$ and $|01 = 101$. In a wing the words face the centre ($4 = 10|,00|,|01$), and a whole number is written to its centre and folded: $8 = 10|,,00|,,10|,,00|$ unfolds to $101,,000,,101,,000,,101,,000,,101$. The half-words $|00$, $|10$, $|01$, $|11$ give back the four words $000$, $010$, $101$, $111$.

Commas count the zoom. One comma separates words, two separate blocks one level out, three the next, so the number of commas at a boundary is the number of zoom levels that end there:
\[ \scriptstyle 17 = 8,,,9,,,8 = 101,000,101\,,,\,000,000,000\,,,\,101,000,101\;,,,\;101,000,101\,,,\,101,111,101\,,,\,101,000,101\;,,,\;\cdots \]
The same string is read as a microscope or a telescope by choosing a comma level. A mark standing alone between commas is a fold at that level: $,|,$ folds at the word level and $,,|,,$ at the block level, and the fold shares the last piece written before it. So $4 = 101,000,|$ (words $101,000$ folded on $000$) and $8 = 4,,0,,|$ (blocks $4,0$ folded on the zero block); the shorthand carries the level at which each mirror acts. The width steps up exactly after the numbers $3\cdot2^k$: $4$--$6$ are $9$ digits wide, $7$--$12$ are $27$, $13$--$24$ are $81$, $25$--$48$ are $243$.

A mark may not take the seat of a unit. A centre that is a zero holds its mark in the zero ($x|x$, the mark in the place of absence); a centre that is a unit keeps it, bracketed ($x|1|x$). With this rule the marked string unfolds to the build exactly (checked for $n<400$), and
\[ \scriptstyle 11 = 5,6,5 = 10|,1|1|1,|01\;,,\;11|,0|0,|11\;,,\;10|,1|1|1,|01. \]
The centre digit of a whole number is a unit exactly for $1$, $3$ and $2^m+1$ (proven in Lean for every $n$, \texttt{centre\_unit\_iff}).

\paragraph{Cube form.} Written with a shared centre, $a,|a|,a$ is $a\times a\times a$: $2,|2|,2 = 8$, $3,|3|,3 = 27$, $4,|4|,4 = 64$. With bare units the counting rule gives $1,|1|,1 = 1\times1\times1 = 3$, so three is the one of cubes, as two is the one of squares ($1\times1 = 2$), and $n$ is the one of $n$-th powers.

\section{The ladder}
At every level of the build of $n$ only two neighbours $k$ and $k+1$ occur, so a number of any size is written by one pair per level, from the top down:
\[ 19 = 9,10,9 \;\big|\; 9|10 : 4|5 : 2|3, \] \[ 97 = 48,49,48 \;\big|\; 48|49 : 24|25 : 12|13 : 6|7 : 3|4 : 1|2. \]
Every ladder ends in $2|3$ (for $n$ from $2^m$ up to $3\cdot2^{m-1}$) or in $1|2$ (from $3\cdot2^{m-1}$ up to $2^{m+1}$), for every $m\ge2$, and the ladder of $n\ge3$ has $\lfloor\log_2\lfloor n/3\rfloor\rfloor+1$ levels (proven, \texttt{ladder\_ends\_two\_three}, \texttt{ladder\_ends\_one\_two}, \texttt{ladder\_length}; the last found by PySR). The number of units climbs the same ladder:
\[ u(1),u(2),u(3) = 1,2,3,\qquad u(2k) = 2u(k),\qquad u(2k{+}1) = 2u(k)+u(k{+}1), \]
so $19$ holds $(2,3)\to(4,7)\to(15,14)\to 44$ units. The sequence $u$, which runs 1, 2, 3, 4, 7, 6, 10, 8, 15, 14, 20, 12, 22, 20, 28, 16, \dots, has the same $k$, $k+1$ shape as Stern's diatomic sequence ($a(2n)=a(n)$, $a(2n{+}1)=a(n)+a(n{+}1)$), a different count reached by a different route; $u$ itself was not found in the literature or in the OEIS.

\section{The ratio-two line}
Since $1\times1 = 2$, two units per number is the natural line to measure against.

\textbf{Theorem (proved in Lean 4 with Mathlib).}
\begin{enumerate}
\item $u(n) = n$ exactly when $n = 2^a$ or $n = 3\cdot2^a$; every other $n$ has $5u(n)\ge 7n$.
\item $u(n)\ne 2n$ for every $n\ge1$: no number sits on the line.
\item $u(2^a+1) = 2^{a+1}-1$ for $a\ge1$, and $u(3\cdot2^a+1) = 2(3\cdot2^a-1)$ for $a\ge0$.
\item An odd number $n$ has $u(n) < 2n$ if and only if $n = 2^a+1$ or $n = 3\cdot2^a+1$ ($a\ge1$), or $n\in\{1, 11, 15\}$.
\end{enumerate}
So below the line there are two endless strands, which fall short by exactly $3$ and exactly $4$, and three single points. The points $11$ and $15$ are forced: an odd number can dip only if one of its blocks counts exactly itself, and if only the middle block does, both halves inside it must as well, which happens only at $11 = 5,6,5$ (wings $2^2+1$ around $3\cdot2$) and $15 = 7,8,7$ (wings $3\cdot2+1$ around $2^3$): the same parts with the roles exchanged. Numerically, $u(n) = 2(n-1)$ holds at $n = 2$, $11$, $15$ and at no other $n<2^{20}$: $11$ holds $20 = 2\cdot10$ and $15$ holds $28 = 2\cdot14$. Up to $2^{20}$ the decimal mirrors below the line are $11$, $22$, $33$, $44$, $66$, $88$, $272$, $2112$, $4224$, $8448$ (checked for $n<2^{20}$), and $u(n)/n$ is exactly $20/11$ on $11\cdot2^j$ and $21/11$ on $33\cdot2^j$. The two strands are also exactly the odd numbers whose count is not a multiple of $4$: for odd $n$, $u(n)$ is odd exactly at $1$ and $2^a+1$, is $2$ more than a multiple of $4$ exactly at $3\cdot2^a+1$, and every other odd $n$ has $4\mid u(n)$ (proven in Lean for every $n$, theorems \texttt{u\_odd\_iff}, \texttt{u\_two\_mod\_four\_iff}); $11$ and $15$ are the first odd numbers whose count holds exactly one factor $4$. The factor comes from the step $2u(n)-u(n{+}1)$, which is $4u(k)$ at $n = 2k+1$.

\section{Even primes and midpoints}
On the even line $2$ is the unit: $2k$ stands at place $k$, and the product of places is $a\otimes b = 2ab$, so $1\otimes1 = 2$, which is $2\times2 = 4$; the counting rule appears one level down. ``Even prime'' then reads two ways: $2p$ (a prime place), or an even number that is not a product of two evens ($n\equiv 2 \pmod 4$, $n\ne2$). In the facing build the second kind is exactly the evens whose code opens with one hinge mirror before an odd core. The two readings first part at $18 = 2\cdot9$. Inside the evens a number may split two ways into such primes: $36 = 2\cdot18 = 6\cdot6$, $60 = 2\cdot30 = 6\cdot10$, $180 = 6\cdot30 = 10\cdot18$.

An odd prime $p = 2k+1$ is $k,\,k{+}1,\,k$, and its midpoint block $k+1$ is even exactly when $p\equiv 3\pmod 4$. The primes with an odd midpoint, $p\equiv1\pmod 4$, are exactly the odd primes that lie on a circle $x^2+y^2 = p$ (the same set appears, by a different route, in Fermat's and Euler's work on sums of two squares).

\section{Two mirrors}
A triple $a,b,a$ can be read two ways. As a \emph{reflection} (the halving mirror, two-based) the far wing is the image of the near one: a hinge reads wing $+$ wing ($8 = 4,0,4$) and a filled centre reads centre $+$ one wing ($7 = 3,4,3$, $4+3$). Read this way, the halving build gives exactly $n$ for every $n$. As a \emph{turn} (the thirding mirror, a triangle of three equal copies) every copy counts: $27 = 9,9,9$, $15 = 5,5,5$, $3 = 1,1,1$; the ternary build $k,k,k$ / $k,k{+}1,k$ / $k{+}1,k,k{+}1$ reads exactly $n$ this way and its digits hold exactly $n$ units (both checked for $n<3000$). The same triple read by the wrong mirror is the paradox: $9,9,9$ as a reflection is $18$, a $9$ where the hinge should be; $1,1,1$ as a reflection is $2$.

\paragraph{The imbalance.} A filled-centre triple $a,b,a$ reads $a+b$ as a reflection and $2a+b$ as a turn: the two differ by one wing, $a$; hinge triples read alike. Expanding a number by halving and adding the wings of all its filled-centre triples gives exactly $u(n)-n$, the units the halving build carries beyond $n$ (proven in Lean for every $n$, theorem \texttt{u\_eq\_add\_imb}). So the halving drift is the mirror wings counted twice, and the self-count numbers $2^a$ and $3\cdot2^a$ are the numbers with no imbalance. For $22 = 11,0,11$ the top is balanced, and the triples $5,6,5$ and $2,3,2$ inside carry $5+2+2+5+2+2 = 18 = 40-22$.

\paragraph{Three is where they join.} By halving alone, $3 = 1,2,1 = 010,101,010$ (half form $010,10|$). The word $111$ used so far is a turn. With $3 = 1,2,1$ the units satisfy $v(1),v(2) = 1,2$ and the same recursion, and then $v(n) = n$ exactly at $2^a$, $v(n)\ne 2n$, and the odd numbers below $2n$ are exactly $1$, $7$ and $2^a+1$ (proven in Lean for every $n$). The second strand $3\cdot2^a+1$ and the points $11$ and $15$ appear only when $3$ is the turn $1,1,1$: the halving strand belongs to the reflection, and the wall strand with its two forced points is the thirding mirror showing through.

\paragraph{Marks.} Two notations both decode exactly. With $|$ as a mark (wings facing in, $L|,C,|R$, for a reflection; every copy facing on, $X|,Y|,Z|$, for a turn; the comma count giving the level) the strings decode back to the build for $n<3000$, and the thirding form is the shorter ($27$: $49$ symbols against $189$); the turn mark string of $3^k$ has exactly $2\cdot3^k-k-2$ symbols (found by PySR, proven, \texttt{markT\_length}). Without an extra symbol, written one order up, every side seat of a triple is empty, so a $1$ in a side seat can only be a mark ($110$, $011$); this costs three times the length.

\section{One picture, three readings}
The two mirrors fold into one term. Write a picture $a,b,a$ once with a scale $s$ after a colon; it reads
\[ a,b,a\!:\!s \;=\; b + (s-1)\,a, \]
with the halving convention that an empty centre is the mirror showing the wing ($a,0,a\!:\!2 = 2a$). Scale $2$ is the reflection, scale $3$ the turn, and every halving form and every thirding form reads exactly $n$ this way (checked for $n<5000$). Scale $1$ is the third reading: the centre alone, the seat of the observer. One picture read at scales $1,2,3$ steps by one wing:
\[ 1,1,1 \to 1,\,2,\,3,\qquad 2,3,2 \to 3,\,5,\,7,\qquad 9,9,9 \to 9,\,18,\,27, \]
so $1,1,1$ holds the first three numbers and $2,3,2$ the first three odd primes. Since one of $b$, $b+2$, $b+4$ is a multiple of $3$ (three steps of $2$ make a full turn of the $3$ mirror), a picture with wing $2$ gives three primes at $2,3,2$ and, since the multiple of $3$ is then $3$ itself, at no other wing-$2$ picture. The numbers whose halving and thirding forms share their wings are exactly $2$, $3$ and $5$ (proven in Lean, theorem \texttt{shared\_wings}): $2 = 1,0,1$ is the same picture both ways (two observers), $3 = 1,2,1\!:\!2 = 1,1,1\!:\!3$ (the first operator), $5 = 2,3,2\!:\!2 = 2,1,2\!:\!3$; from $7$ on the wings part.

\paragraph{Arrows.} An arrow pointing in marks an operator acting on the centre; an arrow pointing out marks an observer. So $3 = 1{>},1,{<}1$ (operators from past and future on the now), $1,{<}1{>},1$ is the centre looking both ways, and $1{>},{<}1{>},{<}1$ shows the whole triad. The shorthand $a{<}b{>}$ means $a,b,a$: $2 = 1{<}0{>}$, $3 = 1{<}1{>}$, $7 = 3{<}1{>}$, and nesting $k \mapsto k{<}1{>} = 2k+1$ runs through $1, 3, 7, 15, 31, \dots$

\section{Every prime a mirror}
Let each number take its first mirror: the least $q$ that splits it into $q$ equal copies. For $q=2$ this is the hinge ${<}k,0,k{>}$; for odd $q$ it is $k{>},{<}k{>},\dots,{<}k$, the operators at the two ends and observers between; $25$ is five $5$s and $35$ five $7$s. A number with no equal split is a prime, and it can only take the halving $k{>},k{+}1,{<}k$: operators acting on a centre that is not a copy of them. Stopping at the mirrors $2$ and $3$ leaves the first false prime at $25$; each added mirror moves it to the next prime squared ($49$, $121$, \dots). The mirrors and their reading are the author's.

A prime's centre $(p+1)/2$ meets its own first mirror. For the primes from $5$ to $10^6$ that mirror is $2$ for $39321$ primes, $3$ for $19610$, $5$ for $4868$, $7$ for $2478$, and for $4299$ primes the centre is itself a prime. Stepping in while the centre stays prime gives strands made only of operators, $q \to 2q-1$; the longest below $10^6$ is
\[ 16651 \to 33301 \to 66601 \to 133201 \to 266401 \to 532801, \]
and most strands are cut after one or two steps. The same prime sequences appear, reached by a different route, as Cunningham chains of the second kind; reading them as strands of pure operators is the author's.

\section{One code, levels down}
The forms above join into a single code. Every number is written $w,c,w$ with $w$ its half (rounded down) and centre $c$ equal to $0$ (the hinge, for an even number) or $1$ (the unit, for an odd one); a comma marks one level down, and each block is written again as its own $w,c,w$, with $0\to0,0,0$ and $1\to0,1,0$ at the bottom:
\[ 8 = 4,0,4 \;\to\; 2,0,2,,0,0,0,,2,0,2,\qquad 7 = 3,1,3 \;\to\; 1,1,1,,0,1,0,,1,1,1,\qquad 11 = 5,1,5. \]
At the bottom the words are $000$, $010$, $101$, $111$, and every number holds exactly itself in units. Equal copies are plain commas ($2,2,2 = 6$), a shared centre marks a product ($2,|2|,2 = 8$, $1,|1|,1 = 3$), and $4,0,4 = 8$ is the halving, all in one notation.

\section{Missing wings and the imaginary side}
An odd prime is written by its wing, $p = |w$: $3 = |1$, $5 = |2$, $7 = |3$, $11 = |5$, $13 = |6$, $17 = |8$, $19 = |9$. A wing goes missing when it is made of copies of an odd block $m$ plus half of $m$, $w = a\,m + (m-1)/2$; each odd block lays a strand that starts at $m + (m-1)/2$ and steps by $m$ ($3$: $4,7,10,13,\dots$; $5$: $7,12,17,\dots$; $7$: $10,17,24,\dots$). The wings on no strand are the prime seats, and laying the strands across a short window ahead of a prime names the next prime as the first open seat. This is proven in Lean for every wing: a wing $w\ge1$ is missing exactly when $2w+1$ is not prime (\texttt{Upgrades.open\_iff\_prime}), and $4w+2$ is twice a prime exactly when $w$ is open (\texttt{twice\_prime\_iff\_open}). Up to a million wings, $148932$ are open; $74416$ of the open wings are even and $74516$ odd. Every even open wing belongs to a prime that is a sum of two squares (proven in Lean for every wing, \texttt{even\_wing\_iff\_two\_squares}) and splits into a mirror pair $a\pm bi$ ($5 = (2+i)(2-i)$, $13 = (3+2i)(3-2i)$), and no odd one does: the even wings are the imaginary side.

\section{The mirror at the walker}
Write the gap between consecutive primes $p < q$ as the triple $h,m,h$: half gap, centre, half gap ($2,9,2$ for $7|11$; $3,26,3$ for $23|29$). The centre's seat fixes the half gap's seat: if $6\mid m$ then $h$ is prime to $6$; if $m$ is an odd multiple of $3$, $h$ is even and prime to $3$; if $m$ is even and prime to $3$, $h$ is an odd multiple of $3$; if $m$ is prime to $6$, $6\mid h$ (in every gap up to $10^6$). In the code $2k$ and $2k+1$ share the wing $k$, so the two sides of a gap are not alike: the next prime finishes a wing pair whose even half sits inside the gap, while the prime below closes a pair outside it. The right side is broken, and that is the tell.

Walking up from $p$, stand at each point $x$, reflect the laid strands through $x$, and read whether the mirror image $2x-p$ is an open seat; the first such $x$ is the centre and $2x-p$ is $q$. This named $q$ exactly, stopping at the centre, in all $78495$ gaps with $5\le p<q\le10^6$, laying half of each gap. Blocks cross the mirror only through the centre: of $194071$ mirrored missing-seat pairs none carry the same block set, and every block shared by the two sides ($46084$ pairs) is a block of the centre. Knowing only the blocks up to $31$ and those crossing the laid window, the same reading names $q$ in $78.3\%$, $60.3\%$ and $49.7\%$ of gaps for $p$ in $[10^3,10^4)$, $[10^4,10^5)$, $[10^5,10^6)$, against $31.3\%$, $24.3\%$, $20.1\%$ for a guess by the $6$-facing; its misses fall on a seat just before $q$, never past it.

\section{Even primes}
On the even line (place $k$ at $2k$, $2$ as its one, product $a\otimes b = 2ab$, so $2\times2 = 4$ is $1_u\times1_u = 2_u$) primality and factoring trade places. The primes inside the evens are the odd places, $2,6,10,14,18,\dots$, a fixed quarter of the numbers ($\lfloor (X+2)/4\rfloor - 1$ up to $X$, found by PySR and proven for every $X$, \texttt{even\_world\_count}), and they do not thin; but an even number may be built from them in several ways ($36 = 2\times18 = 6\times6$; first with $2,3,4,5$ ways: $36$, $180$, $360$, $840$), and the leaders in number of builds up to $10^6$ are largely divisor-record numbers ($36, 180, 360, 840, 2520, 5040, 7560, 25200, 27720, 45360, 55440, \dots, 720720$). Every wing $w$ gives one such seat, $4w+2$, and it is twice a prime exactly when $w$ is open (every wing to $10^6$); the strands on a missing seat appear one level up as extra builds of its double ($18\to36 = 2\times18 = 6\times6$, $22\to44 = 2\times22$ only). The even line has no unit of its own: its unit is $1$, on the odd line. With $2$ taken as the unit ($a\cdot b = ab/2$) the even primes fall to the doubled primes and factoring is single again, but then $2\cdot2 = 2$ and the counting rule is lost: the doubling in $1\times1 = 2$ is what makes the even primes.

Each doubled prime $2p$ ($p$ odd) sits between two left-out evens, one shallow ($4\times$ odd) and one deep (a multiple of $8$). The deep side is on the left exactly when $p\equiv1\pmod4$ and on the right when $p\equiv3\pmod4$, and its depth is one more than the number of halvings of $p-1$ (or $p+1$) (proven in Lean for every odd $p\ge5$, \texttt{depth\_one\_mod\_four}, \texttt{depth\_three\_mod\_four}): $8\,|\,10\,|\,12$, $12\,|\,14\,|\,16$, $32\,|\,34\,|\,36$, $60\,|\,62\,|\,64$.

\section{Three faces in one code}
A picture keeps one notation and its marks say how to read it: commas add ($w,c,w$ and equal copies $m,m,m$), a centre between bars multiplies ($a,|b|,c = a\times b\times c$), and wings facing in mark a reflection, read as centre plus one wing ($|k,(k{+}1)|,|k$; a hinge reads wing plus wing). Every face reads back to $n$ for $n<3000$:
\[ 3 = 1,1,1 = 1,|1|,1 = |1,2|,|1,\qquad 18 = 9,0,9 = 6,6,6 = 2,|3|,3 = |9,0|,|9. \]
By the counting rule $1\times1\times1 = 3$, so $1$ has no product face, and a prime's product face is the one around the unit, $1,|p|,1$.

\section{Halving shows the neighbours, the turn shows the number}
For odd $n$ the halving picture is wing $(n-1)/2$, centre $(n+1)/2$: it is built from the even neighbours $n\pm1$, the left-outs, which share no block with $n$. So $55 = 27,28,27$ shows $54 = 2\cdot3^3$ and $56 = 2^3\cdot7$, while its own blocks live in its turn faces: five copies of $11$, or $5,|11|,1$. Read as a square difference, centre$^2-$wing$^2 = n$ is the face $1\times n$, and every face $a\times b$ is a picture with centre $(a+b)/2$ and wing $(b-a)/2$; the pictures of $n$ are its faces, one to one, and a prime has one picture, its halving (proven in Lean). The same fact is reached by a different route in Fermat's square-difference reading.

\paragraph{The mark $|\text{\textasciicircum}$.} Let $\text{\textasciicircum}$ take a block one scale step up and $|$ mirror it across that step, keeping both: $|\text{\textasciicircum}x = x(b+1)$ in base $b$. In base $10$, $|\text{\textasciicircum}1 = 11$, $|\text{\textasciicircum}\text{\textasciicircum}1 = 101$, $|\text{\textasciicircum}5 = 55$; in base $2$, $|\text{\textasciicircum}1 = 3$ and $|\text{\textasciicircum}\cdots\text{\textasciicircum}1$ runs $3, 5, 9, 17, 33$, the strand $2^a+1$. So $3$ and $11$ are one symbol at two scales, $22 = 11,0,11$ and $23 = 11,1,11$ are the pictures of $2$ and $3$ with the one of tens as wings, and $(k$ ones$)\times(10^{k-1}+1) = (k-1$ ones$)\,2\,(k-1$ ones$)$: $11\times11 = 121$, $11111\times10001 = 111121111$. The pair $7, 5$ sits on the side-and-diagonal ladder of $\sqrt2$: $7^2 = 2\cdot5^2 - 1$, with $(3,2), (17,12), (41,29), \dots$ (the same fact by a different route as the Pell numbers).

\section{Turns, the hierarchy, and the hop}
\textbf{Theorem (Lean, every set of turns).} With any set of odd turns written as equal copies and halving elsewhere, $u(n)\ne 2n$; the self-counting numbers are $c\cdot2^a$ with $c$ one or a turn; every odd non-turn number below the line is a strand point, a pair point $4j+3$ (with $u = 2n-2$) or a point next to a turn; nothing not self-counting falls below $\tfrac34$ of its units. A turn $t = 2^b+1$ makes the pair point $4t-1$, a turn $t = 2^b-1$ makes $4t+3$; $3 = 2+1 = 4-1$ is both and makes $11$ and $15$.

\textbf{Hierarchy (Lean).} A candidate $p+2h$ falls on the seat of a block $r$ exactly when $h \equiv -p\cdot2^{-1} \pmod r$ ($5$: $h\equiv2p$; $7$: $h\equiv3p$). Seat rules stack in size order $2,3 < 5 < 7 < 11 < 13$; block $r$ sits above the hand-off from the seat before its two seats in each run of $6r$, and takes $1/(r-1)$ of what the smaller blocks leave ($0.2503$, $0.1250$, $0.0626$ measured to $10^6$). Over each run of $6r$ the share is exact: of the $2(r-1)$ seats there off $2$, $3$ and $r$, the hand-off's first seat lands on a seat of $r$ from exactly $2$ (found by PySR, proven for every prime block $r\ge5$, \texttt{above\_share}), and stacked on the blocks below it the share is exactly $6/36$ for $7$ and $30/300$ for $11$ (proven). Seats closed twice stay closed once; the offset-$11$ mirror, the per-block mirror and the lake change the picture and never the verdict.

\textbf{The hop (Lean).} Every prime past $3$ sits on the first or last seat of its block of $6$; a hop to the next prime crosses $w$ walls and takes a seat step $s\in\{-1,0,1\}$, with half gap $h = 3w+2s$; the side mod $4$ flips exactly when $w$ is odd and the facing flips exactly when $s\neq0$. The wall-free hop of $4$ leaves $6k+1$ and lands on $6k-1$, so read from one end it belongs to one family and from the other end to the other. Read both ways the flips balance: $22837$ against $22838$ facing flips and $22572$ against $22571$ side flips up to $10^6$, and on any walk the two counts differ by at most one.

\paragraph{The range.} Stacking seat rules and the strands of passed left-outs closes the range for the next prime; the seats needed for $90\%$ / $99\%$ near $10^6$:
\begin{center}\begin{tabular}{lcccccc}\toprule
maps & walk & mirror at walker & $2,3$ & $+5$ & to $13$ & $95$ packed places\\\midrule
seats & $28/50$ & $14/25$ & $9/17$ & $7/13$ & $5/9$ & $1/2$\\\bottomrule
\end{tabular}\end{center}
The chance that the next open seat holds the prime is nearly flat ($0.39$ to $0.45$), so the nearest open seat stays the likeliest and the range closes by adding maps, not by reordering; the residue that remains is the qQq.

\section{Offsets of eleven}
On the mirror line with offset $11$ ($n_m = 11 - n_s$) the odd primes face even numbers: $5\leftrightarrow6$, $7\leftrightarrow4$, $3\leftrightarrow8$, with $2\leftrightarrow9$, $1\leftrightarrow10$ and $11\leftrightarrow0$. A second offset of $22$ ($22_m = 0_s$, $22_s = 0_m$) faces odd numbers with odd numbers. The odd multiples $11, 33, 55, \dots$ pair each prime with an even number, often an even prime ($33$: $7\leftrightarrow26$, $11\leftrightarrow22$, $19\leftrightarrow14$, $23\leftrightarrow10$, $29\leftrightarrow4$), and the even multiples $22, 44, 66, \dots$ pair primes with primes ($22$: $3\leftrightarrow19$, $5\leftrightarrow17$, $11\leftrightarrow11$; $66$: six pairs). In the first blocks of the per-block mirror the primes face even primes and the odd composites face the evens that are not doubled primes ($9\to2$, $15\to18$, $25\to30$, $35\to42$); the first crossings are $33\leftrightarrow22$, where both carry the mirror's own $11$, and $73\leftrightarrow70$. The author deduced these pairings within this system, without prior knowledge of earlier work; for reference only, related statements in the literature are associated with E.~Lemoine (an odd number as a prime plus twice a prime) and C.~Goldbach (an even number as two primes). Here both stand on one mirror line as neighbouring offsets of $11$.

\section{Nine, cubes and eighteen}
Squared and cubed, the mark reads the rows of the triangle in its own base: $(|\text{\textasciicircum}1)^2 = 1,2,1$ and $(|\text{\textasciicircum}1)^3 = 1,3,3,1$, so $9 = 3^2$ is the picture $1,2,1$ written in base $2$, and $121 = 11^2$, $1331 = 11^3$ in base $10$. The two mirrors of decimal reversal are $11$ and $9$: a two-digit number plus its reversal is $11$ times its digit sum and minus it is $9$ times its digit difference; on the nines reversal is the eleven line, $\mathrm{rev}(9k) = 9(11-k)$, so $18\leftrightarrow81$, $27\leftrightarrow72$, $36\leftrightarrow63$, $45\leftrightarrow54$: $18 = 2\times9$ read backwards is $81 = 9\times9$. A cube $p^3$ carries its root's halving scaled ($27 = 3\times(1,2,1) = 6^2-3^2$) and its own halving reads the turn-neighbours $p^2\pm p+1$ ($27 = 13,14,13 = 1\cdot13, 2\cdot7, 1\cdot13$). The squares of $2, 3, 5, 7$ sit one past a self-counting number $2^a$ or $3\cdot2^a$ ($4 = 3+1$, $9 = 2^3+1$, $25 = 3\cdot2^3+1$, $49 = 3\cdot2^4+1$), and, proven for every prime, no other prime's square does (\texttt{square\_past\_self\_counting\_iff}). Sums of two cubes are composite except $1^3+1^3 = 2$; the first two are $2$ and $9$ (\texttt{cube\_sum\_prime\_iff}). All proven in Lean. Same facts reached by different routes: Nicomachus (cubes), Ramanujan ($1729 = 12^3+1 = 10^3+9^3$).

\section{Mapped out to ten million}
From each prime $p$, pack the open seats already found (the primes up to $p$) from the small end into places that fit under $p$, lay their strands forward, and name the first seat past $p$ that no strand reaches. Every prime from $2$ to $10^7$ ($664579$ primes) is named exactly in turn, using at most $213$ packed places; one place reaches $23$, two reach $167$, ten reach $9403$, $29$ reach $100483$, $78$ reach $1018057$. Exactness once enough blocks are laid is a fact reached by a different route in the classical theory; the construction and its count are the author's.

\section{Two as the zero}
Zero halves into $-1,0,1$ and two is $1,0,1$: the same picture, with the wings facing (they cancel) or counted (they make two). Two is zero with its wings counted, and it is the prime with the hinge at its centre; every other prime has the unit there (proven in Lean). Measured from two's mirror every odd prime reads $1$, so two sets the shape of the line and carries no prime information by itself; that begins at the first stacked zero, $3$. Two is $0$ on the number line and $1$ on the even line. Every prime is the zero of its own mirror, and the walls $6, 30, 210, \dots$ are where zeros stack. Around a stack $W$ the open seats pair $r\leftrightarrow W-r$ with the same parity and none its own mirror; within half a turn of $6$, $30$ and $210$ every open seat is $1$ or a prime, which holds while $W/2$ is below the next block squared and first frays at $2310$, at $169 = 13\times13$ (all proven in Lean). Twins straddle a multiple of $6$; for a $6k-1$ prime $p-2$ lands on three's zero and $p+2$ does not, and for a $6k+1$ prime the reverse (to $10^7$).

\section{The wall law}
Every prime past $3$ is a wing of a wall, $p = 6k+f$ with facing $f = \pm1$. A prime block $r = 6k_r + f_r$ is itself a wing of its own wall, and it closes the wing $f$ of wall $k$ exactly when
\[ k \equiv f\, f_r\, k_r \pmod r: \]
each block reads out its own wall index, signed by the product of the two facings, and one residue of $k$ answers for both wings of a wall (proven in Lean; checked for every block from $5$ to $281$ and every wall to $200000$). The law holds at every width: for walls of $W$ and a block $r = W k_r + f_r$ with $f_r g_r \equiv 1 \pmod r$, $r$ divides $Wk+f$ exactly when $r$ divides $k - f g_r k_r$, with no further condition (\texttt{FableImprovements.wall\_law\_width}); walls of $6$ are the case $g_r = f_r$. The two families are one strand through zero, $6(-k)+1 = -(6k-1)$, and a hop between consecutive primes is $2h = 6w + (f_q - f_p)$ with $w$ walls crossed (proven). Read from a prime $p$ alone, the next prime is a wing of a wall above $p$ whose index escapes every laid block's signed index; its facing flips in about $57\%$ of hops; its side mod $4$ is kept in $47.5\%$ of hops from $6k+1$ and $38\%$ from $6k-1$; it sits on the same wall as $p$ in about $20\%$ of hops from $6k-1$ (to $10^7$), and in none from $6k+1$: two consecutive primes past $3$ on one wall are its two wings, $6k-1$ and $6k+1$ (proven, \texttt{same\_wall\_twin}). With the blocks to $97$ laid, the next prime lies within $1/1$, $2/4$, $3/5$, $3/6$ seats ($90\%/99\%$) in the bands $[10^3,10^4)$, $[10^4,10^5)$, $[10^5,10^6)$, $[10^6,10^7)$; working in walls names the same seats in the same order at half the reading cost.

The $6k-1$ family leads the $6k+1$ family in every band to $10^7$ (by $22$, $34$, $189$ to $10^5$, $10^6$, $10^7$). Six opposite readings were tried; at the offset $66$ the prime pairs are equal from both families and the unpaired difference equals the lead exactly. No mirror inside the range balances it: it rides in the qQq. The same lean is known by a different route as Chebyshev's bias.

\section{The two-point law}
Let $v$ be the unit count of pure halving ($3 = 1,2,1$) and $d(c) = 2u(c) - u(c+1)$ the step between neighbours. For every $c\ge2$ and $0\le j\le 2^a$,
\[ u(c\cdot2^a + j) = u(c)\,v(2^a+j) - d(c)\,v(j), \]
so the units of any number come from two readings: the block $c$ it sits in and its offset $j$ read by pure halving. Summed over a block,
\[ 3\sum_{j<2^a} u(c\cdot2^a+j) = u(c)\,(4\cdot5^a - 2^a) - d(c)\,(5^a-2^a), \]
and each step of scale multiplies a block's total by $5$. Both are proven in Lean, and the strands $2^a+1$ and $3\cdot2^a+1$, the class-by-$8$ weights and the octave sum follow from them as corollaries. Both laws were found as exact relations by symbolic regression (PySR) before they were proved. On any walk through seats off $2$ and $3$ (consecutive primes from $5$ among them) the gaps $\equiv 2$ and $\equiv 4 \pmod 6$ differ in number by at most one (proven); their mean sizes differ ($15.66$ against $17.14$ to $10^8$), a lean no smooth gap law yet carries.

\section{What is proven}
Every claim in this paper has been audited and placed in one of five kinds: proven in Lean, checked by computer to a stated range, measured, definition, or reading. Of $141$ rows, $77$ are proven in Lean, in one library (\texttt{MirrorNumbersAll}, $15$ modules, $330$ theorems, standard axioms only, no \texttt{sorry}, built clean). The audit table and the library are in the code bundle. What stays checked to a range is said so where it appears: facts about the original halving string (not yet formalized as a string), the divisor build, finite lists, strand and window checks over primes, and every share and count.

\section{A theoretical note: the method as waves}
In this construction every block $r$ lays a strand of period $r$, a wave, and a prime is a seat where every wave laid so far misses; stacking maps closes the range the way adding waves sharpens an interference pattern. Others have read the primes through waves by other routes, and the closest are physical: superpositions of waves whose zeros fall on the primes (Petersen et al., \emph{Phys. Rev. Lett.} 122, 090201, 2019) and an optical demonstration with gratings of prime periods (Li et al., \emph{Optics Express}, 2020). The dual picture writes the primes as a sum of waves, one per zero of the zeta function (Riemann's explicit formula), whose spacings match the energy levels of heavy nuclei (Montgomery's pair correlation, Dyson's random matrices, Odlyzko's computations), and the primes have been read as an effectively limit-periodic pattern with peaks at rational frequencies (Torquato, Zhang and de Courcy-Ireland, \emph{J. Stat. Mech.}, 2018). We found no serious support for a link between primes and atomic shells. This section is a reading, offered as a direction for physics, not a result.

\section{Related work}
Euclid (\emph{Elements}, Book VII) defines number as a multitude of units and does not count the unit itself as a number; Greek mathematics separates count (\emph{arithmos}) from measure (\emph{megethos}). Von Neumann's construction makes each number contain every number before it. Marty (1934) and Krasner introduced hyperstructures, whose operations may return several values. Terrence Howard has publicly claimed $1\times1 = 2$; he gives no scope for the claim and no containment rule. The containment rule, the whole-expression rule, the two zeros and two ones, the two-valued division and the fold notation here are the author's own. The author learned from, and was inspired by, the work of others, and gratefully notes where the same facts were reached by different routes: Stern's diatomic sequence (M.~A. Stern, \emph{J. reine angew. Math.} 55 (1858) 193--220; Allouche and Shallit, \emph{Theoret. Comput. Sci.} 98 (1992) 163--197; Calkin and Wilf, \emph{Amer. Math. Monthly} 107 (2000) 360--363) has the same $k$, $k+1$ shape with other weights; the set $\{2^a, 3\cdot2^a\}$ is listed as OEIS A029744; the primes $4n+1$ as sums of two squares appear in Fermat (1640) and Euler (E241, \emph{Novi Comm. Acad. Sci. Petrop.} 5, 3--13); the operator strands are the same prime sequences as Cunningham chains of the second kind; the reflected binary code of Stibitz (1941) and Gray (1947) is another mirror built by halving. The method in this paper (the unit count, the two mirrors, the facing marks, the imbalance and the reading by operator and observer) is the author's own. R.~Lipton (blog, 2018) treats $1\times1=2$ as the rescaled product $A\otimes B = 2AB$, which is distributive; the system here keeps containment ($1\times n = n$) and has no distribution.

\section*{Acknowledgements}
The computations, proofs in Lean and drafting were carried out with Anthropic's Claude models: Claude Opus, Claude Fable and Claude Sonnet, working as a team of agents. The author thanks them, and would be glad to be on that team officially. The author also extends a standing invitation to Levent Alp\"oge to join an informal seminar on the rapid exothermic oxidation of certain cannabinoids, loaded into a combustion vessel well beyond its rated capacity; attendance optional, mathematics guaranteed.

\section*{Code and licence}
{\sloppy\texttt{fold\_build\_v2.1\_11-24-2026.py}, the facing-code scripts (\texttt{facing-number-line}, \texttt{ladder-shorthand}, \texttt{eleven-mirror}, \texttt{even-primes}) and the Lean projects \texttt{UnitCountingLaws.lean} and \texttt{RatioTwoLine.lean} (build log and axiom check included; every theorem uses only Lean's standard axioms) implements the counting rule (\texttt{unit\_product}), the halving build (\texttt{B}), orders (\texttt{order}) and the divisor build (\texttt{faces}) with self-checks. Code: MIT. Paper: CC BY 4.0.\par}

\section{Every prime, not a sample}
The laws of the gaps are proven in Lean for every prime past $3$, with no range and no count of places: the hand-off (\texttt{PrimeHandOff.hand\_off\_primes}: for primes $3<p<q$ with $q = p+2h$, the $2$-mirror sits on exactly one of the centre $p+h$ and the half gap $h$, and so does the $3$-mirror), the walker section (\texttt{centre\_fixes\_half\_gap}), no block through both the centre and the half gap (\texttt{never\_both\_primes}), the hop $h = 3w + 2s$ with $s\in\{-1,0,1\}$ and $w$ walls crossed (\texttt{hop\_primes}), and the wall law (\texttt{wall\_law\_prime}, \texttt{wall\_both\_wings\_prime}). They hold at $2^{64}$ and at every size past it.

The walk itself is proven. A seat $n = 6k+f$ beside a wall ($f=\pm1$, $n\ge5$) is prime exactly when no prime block $r = 6k_r+f_r$ with $r\ge5$ and $r\cdot r\le n$ has $k\equiv f f_r k_r \pmod r$ (\texttt{WallWalk.wall\_walk}), and every prime past $3$ is such a seat (\texttt{WallWalk.wall\_walk\_seat}). The blocks a seat needs are the primes up to its square root, which the walk has already named, so the walk from $2$ and $3$ alone names every prime in order, at every size. The runs below show the proof at work; they are not the evidence for it.

Starting from $2$ and $3$ alone, the walk names each prime by the wall law and lays every prime it names as a block for the walls further out; no outside list of primes is used to build anything. Every seat from $2$ to $10^{12}$ was walked on one computer ($12$ cores, $185$ seconds): $37607912018$ primes, equal to the independent count (\texttt{sympy.primepi}) at every power of ten on the way. Every seat from $2$ to $10^{13}$ ($2269$ seconds): $346065536839$ primes, again equal to the independent count. With the blocks the walk named up to $2^{32}$ ($203280275$ of them), the last $1000000002$ seats below $2^{64} = 18446744073709551616$ were built the same way: $67611645$ primes named in order, the last $2^{64}-59$; two separate Miller--Rabin base sets agree on every one, and every one of the $932388357$ closed seats carries a block that divides it. Windows at $10^{12}$, $10^{15}$ and $10^{18}$ gave the same. On walls of $30$, where every prime from $7$ on sits on one of eight wings and the walk is again exact (\texttt{wall\_walk\_thirty}), the same counts come out faster: $10^{11}$ in $2.5$ seconds, $10^{12}$ in $38.1$ seconds and $10^{13}$ in $1051$ seconds, against $16.6$, $185$ and $2269$ for walls of $6$, every count exact. Code and output: \texttt{wall-walk-far}, \texttt{wall-walk-2to64}, \texttt{wall-walk-max}, \texttt{wall-walk-thirty}.

\section{The walls and the Riemann hypothesis}
The Riemann hypothesis has a complete Lean proof (L.~Alp\"oge with Claude; condensed by the Lean community, \texttt{leanprover-community/riemann-hypothesis}), which states Mathlib's own \texttt{RiemannHypothesis} and uses only the standard axioms; the author's group built it from source and checked it. In one Lean environment with that proof, the wall count is tied to it:
\[ \exists\, \delta>0,\ C,\ \forall N\ge3:\quad \Bigl|\, 2 + \#\{\, n\le N : n \text{ passes the wall test} \,\} - \operatorname{Li}(N) \Bigr| \le C\, N^{1-\delta} \]
(\texttt{walls\_meet\_rh}, standard axioms only). The chain runs from \texttt{riemannHypothesis} to a bound on $\zeta'/\zeta$ beside the critical line, through Ingham's transfer theorem as formalized in Tau Ceti, to a power-saving bound on $\pi(N) - \operatorname{Li}(N)$, and then to the wall count by the wall theorem of the section above (the wall count equals $\pi(N)$). Without the Riemann hypothesis the best known error is not a power saving, so this statement about the walls rests on it. The sharper forms are proven too, in the same environment (\texttt{WallsMeetRHSharp}, standard axioms only). Writing $W(N) = 2 + \#\{\, n\le N : n \text{ passes the wall test} \,\}$, which equals $\pi(N)$ by the wall theorem:
\[ \forall \delta>0\ \exists C:\ |W(N) - \operatorname{Li}(N)| \le C\, N^{1/2+\delta}\ (N\ge3), \qquad \exists C':\ |W(N) - \operatorname{Li}(N)| \le C' \sqrt{N}\,(\log N)^2\ (N\ge55) \]
(\texttt{walls\_meet\_rh\_half}, \texttt{walls\_meet\_rh\_sqrt}). They come from the proved hypothesis through a bound on $\zeta'/\zeta$ to the right of the critical line, uniform in the distance to it, and Tau Ceti's smoothed Perron contour estimate; on the way, $\psi(x) - x = O(x^{1/2+\delta})$ and $\psi(x) - x = O(\sqrt{x}\,(\log x)^3)$ are proven. So, because the Riemann hypothesis is true, the walls count the primes to within $\sqrt{N}$ up to a factor $(\log N)^2$. Von Koch's exact $\psi(x) - x = O(\sqrt{x}(\log x)^2)$ and Schoenfeld's explicit constants are not formalized here.

\section{One formula, said plainly}
For readers who want the primes as one formula: write a seat as $n = 6k+f$ with $f=\pm1$, and every block as a seat $6j+g$ with $j\ge1$, $g=\pm1$, named by its own wall $j$ and wing $g$. Then, for $n\ge5$,
\[ n \text{ is prime} \iff \text{for all } j\ge1,\ g=\pm1 \text{ with } (6j+g)^2\le n:\quad (6j+g)\nmid k - f g j \]
(proven in Lean, \texttt{WallFormula.wall\_formula}). It needs no list of primes: the test is one formula in $k$ and $f$. From it the $n$-th prime ($p_0 = 2$) is one count,
\[ p_n + 1 = \#\{\, a \le 2^{n+1} : \pi(a) \le n \,\}, \]
where $\pi(a)$ is the number of primes below $a$, namely $2$, $3$ and the seats below $a$ that pass the test, so that the count of primes up to $N$ is $2$ plus the number of seats that pass, proven equal to Mathlib's own prime count (\texttt{prime\_count\_walls}); the bound $2^{n+1}$ is from the fact that the next prime is at most twice the last (proven, \texttt{nth\_prime\_formula}, \texttt{nth\_prime\_le\_pow}). Said plainly, this formula is a computation. The test is quick; the $n$-th prime formula is exact but adds $2^{n+1}$ terms, so it answers the request for a formula and is not a way to find large primes, which the wall walk does. The author's view stands that the primes are a relationship with families and that the qQq remains.

\section{What is borrowed and what is new}
Every correct way of naming primes ends with the same list, so the list is the check on a method, not the claim of one. This method uses working parts that other methods also use: counting by sixes, remainders, and the fact that a number that is not prime has a block no larger than its square root. The author's account of that use: ``Borrowing working mechanisms from other methods and fixing how they don't work isn't that method. It's seeing where that method is flawed and fixing it.''

What the method does with those parts is the author's, and it is what is claimed here, proven in Lean for every size: each prime past $3$ sits on a wing of a wall of $6$; each block reads the walls it closes from its own wall index, $k\equiv f f_r k_r \pmod r$, one residue for both wings; each split between primes hands its $2$- and $3$-mirrors to its centre and half gap and fixes the half gap from the centre, and for every prime block the centre sits on neither mirror seat of the half gap (\texttt{hand\_off\_every\_block}), the $2$ and $3$ hand-offs being its first two cases; the hop between primes is counted in walls crossed, $h = 3w+2s$; the open wings are the prime seats and the even ones are the primes that are sums of two squares; the walk named by these laws is exact at every size; and beside the primes, the unit count $u$ with its self-counting numbers $2^a$ and $3\cdot2^a$, its ratio-two line that no number touches, and its two-point law. A list of primes gives none of these.

\section*{Version note}
Version 2.0 corrects the fold notation: in version 1.0 the words for $1$ and $2$ were written at order $0$ and the word for $3$ at order $1$ ($010,010,010$). Version 2.0 writes all words at one order, adds orders and the divisor build. Version 2.1 corrects the divisor build: in version 2.0 the first gap between copies was $000$, the same as the zero inside a prime, so products of two primes of $5$ or more ($25$, $35$, $49$, \dots) read as primes; the first gap is now $000,000,000$. Version 3.0 adds the triaxial theorem, the ternary build and the binary drift law, and the mirror, $2$-and-$5$ and carry theorems, all proved in Lean 4. Version 4.0 adds facing marks, the zoom with comma levels, the mark rule, cube form, the ladder shorthand, the ratio-two line theorem (proved in Lean), even primes and midpoints, and the two mirrors (reflection and turn) with the imbalance law, one picture with three readings, arrows for operator and observer, and every prime as a mirror. Version 4.1 rewords the related-work notes: matching known facts are noted as the same facts reached by different routes, not as the source of the method. Version 4.2 adds the Lean proofs of the imbalance law, the multiple-of-four strands and pure halving, the one code with levels down, missing wings and the imaginary side, the mirror at the walker, even primes, three faces in one code, halving as the neighbours and the turn as the number, the mark $|\text{\textasciicircum}$, the Lean results on every set of turns, the hierarchy of seat rules and the hop. Version 5.0 adds an opening statement of what the paper is and is not, offsets of eleven, nine, cubes and eighteen, the construction mapped out to ten million, and a theoretical note on the method as waves. Version 5.1 adds two as the zero and the wall law. Version 5.2 adds the two-point law and the block sum (proven in Lean), the audit of every claim with one combined Lean library, the author's framing of the method, and corrects the no-excess list for the ternary build. Version 5.3 removes wording in versions 4.0 to 5.2 that named another method; that wording was added by the computing assistant against the author's instruction. The strands, walls and wall law are the author's construction. It adds the section on every prime, not a sample, with the wall walk to the edge of $2^{64}$. Version 5.4 adds the Lean proof that the wall walk names exactly the primes at every size; proves in Lean, for every size, claims that versions 5.0 to 5.3 had checked to a range (open wings are prime seats, even open wings and sums of two squares, the centre digit, the zoom width, the ladder ends, the depth law, the squares one past a self-counting number, sums of two cubes); adds four laws found by PySR and then proven (the exact $1/(r-1)$ share over each run of $6r$, the ladder length, the even-line quarter, the turn mark length); renames the numbers $2^a$ and $3\cdot2^a$ in the squares statement as self-counting numbers, keeping wall for a multiple of $6$; walks every seat to $10^{13}$ and the last $10^9$ seats below $2^{64}$; and adds the section on what is borrowed and what is new. Version 5.5 adds the section on one formula, said plainly: the prime test in wall coordinates and the count formula for the $n$-th prime, both proven in Lean (\texttt{WallFormula}). Version 5.6 states the work as a solve of the prime pattern as a structure (where every prime sits, the laws between consecutive primes, and a walk that names each one exactly, proven for every size; not a formula that skips the computation; the qQq remains); adds the wall law at every width with no coprimality clause, the walk on walls of $30$, the wall count proven equal to Mathlib's prime count, one hand-off for every prime block tied to the hop, and the proof that consecutive primes on one wall are its two wings (\texttt{FableImprovements}); and adds the walls-of-$30$ walker. Version 5.7 adds the section on the walls and the Riemann hypothesis (\texttt{walls\_meet\_rh}, checked in one Lean environment with the community proof of the hypothesis) and rewords one history note. Version 5.8 adds the sharp forms, $|W(N)-\operatorname{Li}(N)| \le C N^{1/2+\delta}$ and $\le C\sqrt{N}(\log N)^2$ (\texttt{WallsMeetRHSharp}).
\end{document}
